% Network and survey figures in this section are inputs (OpenStreetMap, JICA tables), listed in
% claims.md under "Inputs taken from documents"; simulated numbers are macros from numbers.tex.
\subsection{Simulator and Network}
The simulation uses UXsim 1.14.2 \citep{seo2023,seo2025}, a link based kinematic wave model
equivalent to a mesoscopic version of Newell's simplified car following model. The network is
imported from OpenStreetMap: every trunk, primary, secondary and tertiary road in a 12 by 12.5~km
box covering Kathmandu and Lalitpur inside and just beyond the Ring Road, 472 nodes and 1,242
directional links, strongly connected. Residential streets, unclassified roads and lanes are not
modeled.

Tertiary roads not in the main road network (194 ways, 391 links) are modeled as side streets:
one lane, a free flow speed of 30~km/h, and an outflow capacity of 0.353 times a main road lane,
the ratio of a narrow two lane road to a two lane road inside the Ring Road in JICA's capacity
table (6,000 against 17,000 PCU per day). Main roads use 60~km/h, JICA's free speed for good
asphalt. A junction where one whole approach group is side streets gets no signal and those
streets give way (merge priority 0.2, an assumed value tested below); the other 181 junctions with
three or more approaches run fixed time signals on a 90 second cycle.

\subsection{Demand}
Demand comes from the JICA 2011 survey \citep{jica2012}. The vehicle origin destination table
gives daily flows between 9 large zones by vehicle type in passenger car units (840,775 PCU per
day); JICA calibrated it to counts at 50 points within about 20 percent. Each large zone pair is
split across the 51 small zones of the person trip table in proportion to person trips, and a peak
hour factor of 0.10 gives 58,525 PCU per hour. Small zones are located by georeferencing the
survey's zone map against geocoded village names (41 control points, median residual 0.59~km) and
load at up to four nearby junctions. The peak is held steady for two hours (\BaseVeh{} PCU in all)
in a four hour horizon. The simulator releases each pair's demand in platoons of five and, through
floating point accumulation, can drop a pair's last platoon; adding half a vehicle to each pair's
requested volume removes the loss, so every run generates exactly the vehicles requested and both
arms of every comparison generate the same number.

\subsection{Calibration and Validation}
Road behavior is calibrated only to JICA's weekday morning travel speed survey (Table 6.2.10):
24.0~km/h on the Ring Road and about 15.3~km/h on inner main routes. Side street capacity and free
flow speeds are taken from JICA's tables rather than fitted. The give way junction rule was added
after signalizing every side street junction held Ring Road speed near 14~km/h. Intersection
counts were read only after each configuration was chosen on speed.

Validation uses the entering peak hour volumes counted by JICA and the Department of Roads at nine
intersections in December 2011 and January 2012, transcribed from the turning movement diagrams in
the survey's appendix, each checked to balance. Model volume is measured over the second hour of
the steady peak, so the empty network at the start does not depress it. Fit is reported as the
model to observed ratio, the correlation across junctions, and the GEH statistic.

\subsection{Interventions}
H1 replaces fixed time signals with a queue actuated heuristic that switches phase when a competing
queue exceeds the current one by a lost time adjusted threshold. H2 tolls each link at 0.2 times
its own delay in an untolled baseline run (\TolledLinks{} links); demand does not respond to price,
so tolls only reroute trips. H3 applies both. H4 compares morning demand (trips toward the core
raised by 40 percent, away from it lowered) with and without one lane shifted toward the peak
direction; on this network two of the three original corridors are found. H5 raises free flow
speeds by the same 35/30 ratio as originally (70 and 35~km/h). H6 spreads the same total demand
over a 50 percent longer window. H7 cuts trips into the ten core zones by a quarter, and is also
compared with a volume matched control that cuts every trip uniformly to the same total.

\subsection{Statistical Design}
Each of the ten configurations runs on the same 15 random seeds, and each test reports the paired
percentage change with a 95 percent confidence interval. The primary measure is mean time in system
over all vehicles, counting unfinished trips up to the horizon, so a lever cannot look better by
stranding slow trips. Mean delay over completed trips, the measure of the first version, is
reported alongside. Eight tests are made; a Bonferroni correction widens each interval by a factor
of about 1.55, and results are labeled by whether they survive it.
